\begin{textsample}{POS Dim 2 – ba – Score 58.00 – ba/ba\_em\_3.txt}  \label{ex:f2_pos_031}
3.
MARCOS LEGAIS > “... é preciso uma aldeia inteira para educar uma criança ”. ( Provérbio africano ) Os marcos legais que determinam as \textbf{políticas} públicas nacionais, inclusive as educacionais, são frutos de uma construção histórica do ordenamento jurídico brasileiro, atravessados pelos acirramentos entre os interesses da população, do Estado e dos organismos internacionais.
Nos países em desenvolvimento, o retardo para a \textbf{efetivação} da educação de \textbf{qualidade} social, em prol da \textbf{equidade} de direitos e de oportunidades para todas as pessoas, se arrasta por gerações, por seu caráter estrutural.
Os marcos legais, nacionais e \textbf{estaduais}, que serão apresentados neste \textbf{capítulo} trazem aspectos considerados relevantes para \textbf{amparar} a organização dos \textbf{currículos} das escolas que \textbf{ofertam} o Ensino Médio no estado, por meio dos seus Projetos Políticos Pedagógicos ( PPP ).
Importante destacar que se faz necessário, no âmbito dos Territórios de Identidade e dos \textbf{municípios} baianos, que as unidades escolares considerem as \textbf{legislações} e normativos educacionais locais, bem como as \textbf{políticas} públicas relacionadas às juventudes, para fomentar a organização curricular dessas escolas.
Vale ressaltar, ainda, que os marcos legais abordados no Volume 1 do DCRB ( Educação Infantil e Ensino Fundamental ) deverão ser considerados na organização curricular do Ensino Médio, bem como todos os Temas Integradores do \textbf{Currículo}, também presentes no referido Volume 1, a saber : Educação em Direitos Humanos, Educação para a Diversidade, Educação para o Trânsito, Saúde na Escola, Educação Ambiental, Educação Financeira e para o Consumo, Cultura Digital e Educação Fiscal.
Essas considerações sobre o Volume 1 são importantes, pois são aspectos que dão sentido de complementaridade e continuidade entre os dois volumes.
Inicialmente, como no Volume 1, o ponto de partida é a Constituição da República Federativa do Brasil de 1988 que, em seu Artigo 205, define a educação como um direito : “ a educação, direito de todos e dever do Estado e da família, será promovida e incentivada com a colaboração da sociedade, visando ao pleno desenvolvimento da pessoa, seu preparo para o exercício da cidadania e sua \textbf{qualificação} para o trabalho ” ( BRASIL, 1988 ).
Este Artigo é basilar para um Estado democrático de direito, pois compromete o Estado, a partir dos seus entes federados, a família na \textbf{garantia} do direito à educação e, também, serão \textbf{responsabilizados} quando esse direito não for assegurado a todas as pessoas, sem distinção de cor, raça, etnia, gênero, orientação sexual, religião, fase geracional, lugar de moradia e/ou condição de deficiência.
O artigo nº 206, também da Constituição Cidadã, discorre sobre os princípios do ensino, que reforçam o compromisso do Estado com os seus cidadãos : I.
Igualdade de condições para o acesso e permanência na escola ; II.
Liberdade de aprender, ensinar, pesquisar e divulgar o pensamento, a arte e o saber ; III.
Pluralismo de ideias e de concepções pedagógicas e coexistência de instituições públicas e privadas de ensino ; IV.
Gratuidade do ensino público em estabelecimentos oficiais. ( BRASIL, 1988 ) Neste sentido, garantir a educação inclusiva e promover a aprendizagem equitativa é um direito fundamental dos/as estudantes.
Por isso, as \textbf{políticas} públicas educacionais necessitam assegurar a \textbf{equidade} das condições ou meios para consolidar o direito de aprender dos/as estudantes.
Deve -se compreender a \textbf{equidade} educacional enquanto princípio educativo fundamental, levando em consideração que um dos papéis da educação é \textbf{atender} a demanda social, garantindo a aplicabilidade dos Direitos Fundamentais expressos na Constituição da República Federativa do Brasil.
Ainda discorrendo sobre o ordenamento jurídico brasileiro, no âmbito da Constituição Federal, só muito recentemente, em 2009, após 21 anos de sancionada a Constituição, por meio da Emenda Constitucional nº 59, que se tornou obrigatória a \textbf{oferta} da educação básica e gratuita dos 4 aos 17 anos, e a todas as pessoas que a ela não tiveram acesso na idade adequada ; o atendimento aos/às estudantes, por meio dos programas suplementares de material didático escolar, transporte, alimentação e assistência à saúde ; a universalização do ensino obrigatório, \textbf{ofertados} em colaboração entre os entes federados ; e a distribuição dos recursos públicos para assegurar prioridade ao atendimento das necessidades do ensino obrigatório, no que se refere à universalização, \textbf{garantia} de padrão de \textbf{qualidade} e \textbf{equidade}, nos termos do plano nacional de educação.
Esses aspectos da Emenda foram centrais para \textbf{viabilizar} e concretizar, efetivamente, o direito à educação básica no país, muito embora foram disponibilizados tardiamente, pois, muito provavelmente, gerações de pessoas foram prejudicadas pela ausência desse direito.
A Emenda ainda alterou o Artigo 214 da Constituição que tratou de aspectos importantes para melhoria no Plano Nacional de Educação, como > definir \textbf{diretrizes}, objetivos, metas e estratégias de implementação para assegurar a manutenção e desenvolvimento do ensino em seus diversos níveis, etapas e modalidades por meio de ações integradas dos poderes públicos das diferentes esferas federativas que conduzam. ( BRASIL, 2014 ) O Estatuto da Criança e do Adolescente ( ECA ), Lei nº 8.069/90, é outra \textbf{legislação} que ratifica o direito à educação de crianças e \textbf{adolescentes}, mais especificamente, no \textbf{Capítulo} IV Do Direito à Educação, à Cultura, ao Esporte e ao Lazer, conforme disposto no Artigo 53 : > A criança e o \textbf{adolescente} têm direito à educação, visando ao pleno desenvolvimento de sua pessoa, preparo para o exercício da cidadania e \textbf{qualificação} para o trabalho, assegurando -se-lhes : > > I. igualdade de condições para o acesso e permanência na escola ; > II. direito de ser respeitado por seus educadores ; > III. direito de contestar critérios avaliativos, podendo recorrer às instâncias escolares superiores ; > IV. direito de organização e participação em \textbf{entidades} estudantis ; > V. acesso à escola pública e gratuita próxima de sua residência. > > Parágrafo único.
É direito dos pais ou responsáveis ter ciência do processo pedagógico, bem como participar da definição das propostas educacionais.
Já nas alíneas do Art. 54, são considerados aspectos importantes que, caso não sejam asseguradas pelas \textbf{políticas} públicas educacionais voltadas para a Educação Básica, impactam sobremaneira no direito à educação dos/as estudantes que estão ou estarão acessando o Ensino Médio : I.
É dever do Estado assegurar à criança e ao \textbf{adolescente} : II. ensino fundamental, obrigatório e gratuito, inclusive para os que a ele não tiveram acesso na idade própria ; III. progressiva extensão da obrigatoriedade e gratuidade ao Ensino Médio ; IV. atendimento educacional especializado aos portadores de deficiência, preferencialmente na \textbf{rede} regular de ensino ; V. atendimento em creche e pré-escola às crianças de zero a cinco anos de idade ; ( Redação dada pela Lei nº 13.306, de 2016 ) VI. acesso aos níveis mais elevados do ensino, da pesquisa e da criação artística, segundo a capacidade de cada um ; VII. \textbf{oferta} de ensino noturno regular, adequado às condições do \textbf{adolescente} \textbf{trabalhador} ; VIII. atendimento no ensino fundamental, através de programas suplementares de material didático-escolar, transporte, alimentação e assistência à saúde. ( BRASIL, 1990 ) Outra lei que dialoga diretamente com as juventudes e precisa ser conhecida e considerada pelas Unidades Escolares é a Lei nº 12.852/2013, que \textbf{institui} o Estatuto das Juventudes.
Desta lei, embora ela também ratifique o direito à educação, faz -se necessário conhecer os princípios desta \textbf{política}, pois são considerados \textbf{estruturantes} para a organização curricular das escolas que \textbf{ofertam} a etapa do Ensino Médio.
Além do mais, os aspectos concernentes ao Direito à Educação já foram tratados no Volume 1 do DCRB.
A seguir, destaca -se o Art. 2º, que trata dos princípios do Estatuto das Juventudes : > O disposto nesta Lei e as \textbf{políticas} públicas de juventude são regidos pelos seguintes princípios : > > I. promoção da autonomia e emancipação dos jovens ; > II. valorização e promoção da participação social e política, de forma direta e por meio de suas representações ; > III. promoção da criatividade e da participação no desenvolvimento do País ; > IV. reconhecimento do jovem como sujeito de direitos universais, geracionais e singulares ; > V. promoção do bem-estar, da experimentação e do desenvolvimento integral do jovem ; > VI. respeito à identidade e à diversidade individual e coletiva da juventude ; > VII. promoção da vida segura, da cultura da paz, da solidariedade e da não discriminação ; e > VIII. valorização do diálogo e convívio do jovem com as demais gerações. > > ( BRASIL, 2013 ) Os/as estudantes que fazem parte do Ensino Médio brasileiro e baiano, bem como os/as do Ensino Fundamental, são formados, também, por pessoas adultas e idosas, precisamente porque o Estado Brasileiro ainda não conseguiu assegurar \textbf{políticas} públicas que assegurem o acesso da população à educação na idade adequada.
Assim, o Artigo 21 do Estatuto do Idoso, Lei nº 10.741/03, visa marcar o direito à educação das pessoas idosas : “ o Poder Público criará oportunidades de acesso do idoso à educação, adequando \textbf{currículos}, metodologias e material didático aos programas educacionais a ele \textbf{destinados} ” ( BRASIL, 2003 ).
Para a organização curricular do Ensino Médio da Bahia, considerar a diversidade étnico-racial, as especificidades das comunidades tradicionais e das populações em situação de itinerância, a exemplo dos povos ciganos, é algo inegociável.
Dessa forma, será assegurada a aplicabilidade das Leis nº 10.639/2003 e nº 11.645/2008 no DCRB, nas habilidades específicas da Bahia, nos objetos de conhecimento e nos componentes curriculares que abordarão a história e cultura indígena, africana e ameríndia e a Educação Étnico-racial, conforme apresentados nos \textbf{capítulos} que tratarão dos Organizadores Curriculares das \textbf{ofertas} de Ensino Médio, tanto na Formação Geral Básica ( FGB ) quanto nos \textbf{Itinerários} \textbf{Formativos}.
Já a Resolução CNE/CEB nº 3, de 16 de maio de 2012, define \textbf{diretrizes} para o atendimento de educação escolar das populações em situação de itinerância, dando destaque à \textbf{garantia} do acesso ( matrícula ), permanência e conclusão dos estudos dessas populações.
O Artigo 2º dispõe que > visando à \textbf{garantia} dos direitos \textbf{socioeducacionais} de crianças, \textbf{adolescentes} e jovens em situação de itinerância, os sistemas de ensino deverão adequar -se às particularidades desses estudantes. ( BAHIA, 2015 ) A Lei de \textbf{Diretrizes} e Bases da Educação Nacional ( LDBEN ), nº 9.394/96, no Artigo 2º, define os princípios gerais e finalidades da educação, comuns a todos os/as estudantes brasileiros/as : > a educação, dever da família e do Estado, inspirada nos princípios de liberdade e nos ideais de solidariedade humana, tem por finalidade o pleno desenvolvimento do educando, seu preparo para o exercício da cidadania e sua \textbf{qualificação} para o trabalho. ( BRASIL, 1996 ) Ainda na LDBEN, no Artigo nº 3º, delineiam -se os princípios basilares para o ensino, em todos os níveis e etapas : I. igualdade de condições para o acesso e permanência na escola ; II. liberdade de aprender, ensinar, pesquisar e divulgar a cultura, o pensamento, a arte e o saber ; III. pluralismo de ideias e de concepções pedagógicas ; IV. respeito à liberdade e apreço à tolerância ; V. coexistência de instituições públicas e privadas de ensino ; VI. gratuidade do ensino público em estabelecimentos oficiais ; VII. valorização do \textbf{profissional} da educação escolar ; VIII. gestão democrática do ensino público, na forma desta Lei e da \textbf{legislação} dos sistemas de ensino ; IX. \textbf{garantia} de padrão de \textbf{qualidade} ; X. valorização da experiência extraescolar ; XI. vinculação entre a educação escolar, o trabalho e as práticas sociais ; XII. consideração com a diversidade étnico-racial ; XIII. \textbf{garantia} do direito à educação e à aprendizagem ao longo da vida. ( BRASIL, 1996 ) Retomando a atenção para o Ensino Médio, o Artigo nº 35 da LDBEN apresenta as finalidades \textbf{precípuas} do Ensino Médio : I. a consolidação e o aprofundamento dos conhecimentos adquiridos no ensino fundamental, possibilitando o \textbf{prosseguimento} de estudos ; II. a preparação básica para o trabalho e a cidadania do educando, para continuar aprendendo, de modo a ser capaz de se adaptar com flexibilidade a novas condições de ocupação ou aperfeiçoamento posteriores ; III. o aprimoramento do educando como pessoa humana, incluindo a formação ética e o desenvolvimento da autonomia intelectual e do pensamento crítico ; IV. a compreensão dos fundamentos científico-tecnológicos dos processos produtivos, relacionando a teoria com a prática, no ensino de cada disciplina. ( BRASIL, 1996 ) No contexto atual das Políticas Educacionais no país, a sanção da Lei nº 13.005/2014, que aprovou o Plano Nacional de Educação ( PNE 2014-2024 ), é mais um marco importante para a educação no país, pois expressa um pacto social firmado, “ por meio de ações integradas dos poderes públicos das diferentes esferas federativas ” ( BRASIL, 2014 ), para a \textbf{garantia} da educação gratuita, pública e com \textbf{qualidade} social para os/as cidadãos/ãs brasileiros/as, pois nele estão expressos os caminhos para as metas e estratégias.
Embora todas as metas se complementem e se retroalimentem, as Metas 3, 7, 11 e 17 têm uma relação mais direta com a etapa para o Ensino Médio e com os desdobramentos na atual \textbf{política} curricular do Ensino Médio.
A Meta 3 trata da universalização, até 2016, do atendimento escolar para toda a população de 15 ( quinze ) a 17 ( dezessete ) anos e elevar, até o final do período de vigência deste PNE, a taxa líquida de matrículas no Ensino Médio para 85 \% ( oitenta e cinco por cento ) ( Brasil, 2014 ).
Observando as estratégias 3.1 e 3.2 da Meta 3, descritas a seguir, são evidenciados aspectos preponderantes para os desdobramentos das outras \textbf{políticas} voltadas à \textbf{política} curricular do Ensino Médio brasileiro, que serão tratadas mais adiante, conforme descrito a seguir : > institucionalizar programa nacional de renovação do Ensino Médio, a fim de incentivar práticas pedagógicas com abordagens interdisciplinares estruturadas pela relação entre teoria e prática, por meio de \textbf{currículos} escolares que organizem, de maneira \textbf{flexível} e diversificada, conteúdos obrigatórios e \textbf{eletivos} articulados em dimensões como ciência, trabalho, linguagens, tecnologia, cultura e esporte, garantindo -se a aquisição de equipamentos e laboratórios, a produção de material didático específico, a formação continuada de professores e a articulação com instituições acadêmicas, esportivas e culturais ( BRASIL, 2014 ). > o Ministério da Educação, em articulação e colaboração com os entes federados e ouvida a sociedade mediante consulta pública nacional, elaborará e encaminhará ao Conselho Nacional de Educação - CNE, até o 2º ( segundo ) ano de vigência deste PNE, proposta de direitos e objetivos de aprendizagem e desenvolvimento para os ( as ) os ( as ) estudantes ( as ) de Ensino Médio, a serem atingidos nos tempos e etapas de organização deste nível de ensino, com vistas a garantir formação básica comum. ( BRASIL, 2014 ) A BNCC também é citada na estratégia 15.6 da meta 15 do PNE e deverá ser considerada nos \textbf{currículos} dos \textbf{cursos} de licenciatura de todo o país, conforme descrito a seguir : > [... ] promover a \textbf{reforma} curricular dos \textbf{cursos} de licenciatura e estimular a renovação pedagógica, de forma a assegurar o foco no aprendizado do ( a ) aluno ( a ), dividindo a \textbf{carga} \textbf{horária} em formação geral, formação na área do saber e didática específica e incorporando as modernas tecnologias de informação e comunicação, em articulação com a base nacional comum dos \textbf{currículos} da Educação Básica, de que tratam as estratégias 2.1, 2.2, 3.2 e 3.3 deste PNE [... ]. ( BRASIL, 2014 ) Sendo assim, implementar a BNCC, por meio dos referenciais curriculares \textbf{estaduais} e/ou \textbf{municipais}, deverá ser considerado na formação inicial dos/as \textbf{profissionais} de educação, nos \textbf{currículos} dos \textbf{cursos} de pedagogia e licenciaturas e não só na formação inicial, mas, também, na formação continuada dos/as \textbf{profissionais} da educação que já estão na ativa.
Dessa forma, faz -se necessário considerar pelos sistemas, \textbf{redes} e instituições de ensino, em todo o país, a constituição de planos e/ou programas de formação.
Após a sanção do PNE, o Ministério da Educação, ainda na gestão do governo da presidenta Dilma Rousseff, deu início ao processo da escrita da BNCC em junho de 2015.
Em setembro de 2015, a primeira versão da BNCC foi à Consulta Pública durante seis meses, disponibilizada à sociedade por meio digital em plataforma específica organizada pelo MEC.
Em junho de 2016, a segunda versão, com as contribuições da primeira versão, retorna à Consulta Pública, realizada por meio de seminários \textbf{estaduais}.
Em fevereiro de 2017, fruto da transição entre o Projeto de Lei nº 6.840/2013 e a Medida Provisória ( MP ) nº 746/2016, foi sancionada a Lei nº 13.415/2017, que alterou a LDBEN e se tornou o marco das \textbf{políticas} públicas contemporâneas, que deram materialidade à \textbf{reforma} do Ensino Médio.
A Lei 13.415/17 altera > as Leis nos 9.394, de 20 de dezembro de 1996, que estabelece as \textbf{diretrizes} e bases da educação nacional, e 11.494, de 20 de junho de 2007, que regulamenta o Fundo de Manutenção e Desenvolvimento da Educação Básica e de Valorização dos \textbf{Profissionais} da Educação, a Consolidação das Leis do Trabalho - CLT, aprovada pelo Decreto-Lei nº 5.452, de 1º de maio de 1943, e o Decreto-Lei nº 236, de 28 de fevereiro de 1967 ; \textbf{revoga} a Lei nº 11.161, de 5 de agosto de 2005 ; e \textbf{institui} a Política de Fomento à Implementação de Escolas de Ensino Médio em Tempo Integral.
Os parágrafos referentes aos Artigos nº 24, 35 -A e 35, da Lei nº 13.415/2017, impactaram diretamente na \textbf{política} curricular nacional e \textbf{estadual}, são eles : >"Art. nº 24................................................................................................................ > > (... ) § 1º A \textbf{carga} \textbf{horária} mínima anual de que trata o inciso I do caput deverá ser ampliada de forma progressiva, no Ensino Médio, para mil e quatrocentas horas, devendo os sistemas de ensino oferecer, no prazo máximo de cinco anos, pelo menos mil horas anuais de \textbf{carga} \textbf{horária}, a partir de 2 de março de 2017. > > § 2º Os sistemas de ensino disporão sobre a \textbf{oferta} de educação de jovens e adultos e de ensino noturno regular, adequado às condições do educando, conforme o inciso VI do art. 4º."> > (... ) Art. 35 -A.
A Base Nacional Comum Curricular definirá direitos e objetivos de aprendizagem do Ensino Médio, conforme diretrizes do Conselho Nacional de Educação, nas seguintes áreas do conhecimento : > > I.
Linguagens e suas Tecnologias ; > II.
Matemática e suas Tecnologias ; > III.
Ciências da Natureza e suas Tecnologias ; > IV.
Ciências Humanas e sociais Aplicadas. > > § 1º A parte diversificada dos \textbf{currículos} de que trata o caput do art. 26, definida em cada sistema de ensino, deverá estar harmonizada à Base Nacional Comum Curricular e ser articulada a partir do contexto histórico, econômico, social, ambiental e cultural. > > § 2º A Base Nacional Comum Curricular referente ao Ensino Médio incluirá obrigatoriamente estudos e práticas de educação física, arte, sociologia e filosofia. > > § 3º O ensino da língua portuguesa e da matemática será obrigatório nos três anos do Ensino Médio, assegurada às comunidades indígenas, também, a utilização das respectivas línguas maternas. > > § 4º Os \textbf{currículos} do Ensino Médio incluirão, obrigatoriamente, o estudo da língua inglesa e poderão \textbf{ofertar} outras línguas estrangeiras, em caráter optativo, preferencialmente o espanhol, de acordo com a disponibilidade de \textbf{oferta}, locais e \textbf{horários} definidos pelos sistemas de ensino. > > § 5º A \textbf{carga} \textbf{horária} \textbf{destinada} ao cumprimento da Base Nacional Comum Curricular não poderá ser superior a mil e oitocentas horas do total da \textbf{carga} \textbf{horária} do Ensino Médio, de acordo com a definição dos sistemas de ensino. > > [... ] Art. 36.
O \textbf{currículo} do Ensino Médio será composto pela Base Nacional Comum Curricular e por \textbf{itinerários} \textbf{formativos}, que deverão ser organizados por meio da \textbf{oferta} de diferentes arranjos curriculares, conforme a relevância para o contexto local e a possibilidade dos sistemas de ensino, a saber : > > I.
Linguagens e suas Tecnologias ; > II.
Matemática e suas Tecnologias ; > III.
Ciências da Natureza e suas Tecnologias ; > IV.
Ciências Humanas e Sociais Aplicadas ; > V.
Formação \textbf{Técnica} e Profissional. > > (... ) § 3º A critério dos sistemas de ensino, poderá ser composto \textbf{itinerário} \textbf{formativo} integrado, que se traduz na composição de componentes curriculares da Base Nacional Comum Curricular - BNCC e dos \textbf{itinerários} \textbf{formativos}, considerando os incisos I a V do caput. ( BRASIL, 2017 ) Por força da Lei nº 13.415/2017, a organização da BNCC, que outrora fora apresentada à sociedade como um documento único que integrava as três etapas da Educação Básica ( Educação Infantil, Ensino Fundamental e Ensino Médio ), após a sanção da Lei nº 13.415/17, a etapa do Ensino Médio foi subtraída, provisoriamente, da versão final do documento, sendo a BNCC da Educação Infantil e Ensino Fundamental aprovada em dezembro de 2017, após perpassar pelas rodadas de audiências públicas promovidas pelo Conselho Nacional de Educação, entre os meses de junho a setembro de 2017.
Durante esse intervalo de tempo, dezembro de 2017 a abril de 2018, a BNCC da etapa do Ensino Médio sofreu alterações para \textbf{atender} a organização curricular estabelecida pela Lei nº 13.415/2017.
Em abril de 2018, o CNE recepcionou a Base do Ensino Médio pelo MEC e, mais uma vez, coordenou consultas públicas em todas as regiões do país, que acarretou, em dezembro de 2018, por meio da Resolução CNE/CP nº 4, de 17 de dezembro de 2018, na versão final da BNCC do Ensino Médio que \textbf{normatizou} : > \textbf{Instituir} a Base Nacional Comum Curricular na Etapa do Ensino Médio ( BNCC-EM ), como etapa final da Educação Básica, nos termos do artigo 35 da LDB, completando o conjunto constituído pela BNCC da Educação Infantil e do Ensino Fundamental, com base na Resolução CNE/CP nº 2/2017, fundamentada no Parecer CNE/CP nº 15/2017. ( BRASIL 2018 ) Anteriormente à homologação da BNCC do Ensino Médio, o CNE \textbf{atualizou} as \textbf{Diretrizes} Curriculares Nacionais do Ensino Médio ( DCNEM ), por meio da Resolução CNE/CP nº 3, de 03 de outubro de 2018.
Como alguns destaques das DCNEM \textbf{atualizadas} temos a ampliação gradativa da \textbf{carga} \textbf{horária} do Ensino Médio, conforme disposto nos parágrafos 2º e 3º do Artigo 17 : > § 2º No ensino médio diurno, a duração mínima é de 3 ( três ) anos, com \textbf{carga} \textbf{horária} mínima total de 2.400 ( duas mil e quatrocentas ) horas, tendo como referência uma \textbf{carga} \textbf{horária} anual de 800 ( oitocentas ) horas, distribuídas em, pelo menos, 200 ( \textbf{duzentos} ) dias de efetivo trabalho escolar, considerando que : > > I. a \textbf{carga} \textbf{horária} total deve ser ampliada para 3.000 ( três mil ) horas até o início do ano letivo de 2022 ; > II. a \textbf{carga} \textbf{horária} anual total deve ser ampliada progressivamente para 1.400 ( um mil e quatrocentas ) horas. > > § 3º No ensino médio noturno, adequado às condições do estudante e respeitados o mínimo de 200 ( \textbf{duzentos} ) dias letivos e 800 ( oitocentas ) horas anuais, a proposta pedagógica deve \textbf{atender}, com \textbf{qualidade}, a sua singularidade, especificando uma organização curricular e metodológica diferenciada, e pode, para garantir a permanência e o êxito destes estudantes, ampliar a duração do \textbf{curso} para mais de 3 ( três ) anos, com menor \textbf{carga} \textbf{horária} diária e anual, garantido o total mínimo de 2.400 ( duas mil e quatrocentas ) horas até 2021 e de 3.000 ( três mil ) horas a partir do ano letivo de 2022.
No Artigo 11 é tratada a Formação Geral Básica, como ela poderá ser \textbf{ofertada} pelos sistemas de ensino, as suas áreas de conhecimento e \textbf{carga} \textbf{horária} total máxima. > (... ) Art. 11.
A formação geral básica é composta por competências e habilidades previstas na Base Nacional Comum Curricular ( BNCC ) e articuladas como um todo indissociável, enriquecidas pelo contexto histórico, econômico, social, ambiental, cultural local, do mundo do trabalho e da prática social, e deverá ser organizada por áreas de conhecimento : > > I.
Linguagens e suas Tecnologias ; > II.
Matemática e suas Tecnologias ; > III.
Ciências da Natureza e suas Tecnologias ; > IV.
Ciências Humanas e Sociais Aplicadas. > > § 1º A organização por áreas do conhecimento implica o \textbf{fortalecimento} das relações entre os saberes e a sua contextualização para apreensão e intervenção na realidade, requerendo planejamento e execução conjugados e cooperativos dos seus professores. > > § 2º O \textbf{currículo} por área de conhecimento deve ser organizado e planejado dentro das áreas de forma interdisciplinar e transdisciplinar. > > § 3º A formação geral básica deve ter \textbf{carga} \textbf{horária} total máxima de 1.800 ( mil e oitocentas ) horas, que garanta os direitos e objetivos de aprendizagem, expressos em competências e habilidades, nos termos da Base Nacional Comum Curricular ( BNCC ). > > (... ) § 7º A critério dos sistemas de ensino, a formação geral básica pode ser contemplada em todos ou em parte dos anos do \textbf{curso} do ensino médio, com exceção dos estudos de língua portuguesa e da matemática que devem ser incluídos em todos os anos escolares. ( BRASIL 2018 ) Sobre a organização da Formação Geral Básica no Estado, os componentes curriculares continuarão sendo trabalhados, pois no processo de transição do \textbf{currículo} do Ensino Médio, se faz necessária apropriação da nova arquitetura curricular.
Contudo, seguindo o disposto no § 2º do Artigo 11 das DCNEM, haverá de se considerar a organização e planejamento curricular, dentro das áreas de forma interdisciplinar e transdisciplinar. ( BRASIL, 2018 ) As DCNEM, ainda, orientam em detalhe os \textbf{Itinerários} \textbf{Formativos}, a serem \textbf{ofertados} pelo Sistema de Ensino, conforme disposto nos Art. 12 : > A partir das áreas do conhecimento e da formação \textbf{técnica} e \textbf{profissional}, os \textbf{itinerários} \textbf{formativos} devem ser organizados, considerando : > > I. linguagens e suas tecnologias : aprofundamento de conhecimentos \textbf{estruturantes} para aplicação de diferentes linguagens em contextos sociais e de trabalho, estruturando arranjos curriculares que permitam estudos em línguas vernáculas, estrangeiras, clássicas e indígenas, Língua Brasileira de Sinais ( LIBRAS ), das artes, design, linguagens digitais, corporeidade, artes cênicas, roteiros, produções literárias, dentre outros, considerando o contexto local e as possibilidades de \textbf{oferta} pelos sistemas de ensino ; > > II. matemática e suas tecnologias : aprofundamento de conhecimentos \textbf{estruturantes} para aplicação de diferentes conceitos matemáticos em contextos sociais e de trabalho, estruturando arranjos curriculares que permitam estudos em resolução de problemas e análises complexas, funcionais e não-lineares, análise de dados estatísticos e probabilidade, geometria e topologia, robótica, automação, inteligência artificial, programação, jogos digitais, sistemas dinâmicos, dentre outros, considerando o contexto local e as possibilidades de \textbf{oferta} pelos sistemas de ensino ; > > III. ciências da natureza e suas tecnologias : aprofundamento de conhecimentos \textbf{estruturantes} para aplicação de diferentes conceitos em contextos sociais e de trabalho, organizando arranjos curriculares que permitam estudos em astronomia, metrologia, física geral, clássica, molecular, quântica e mecânica, instrumentação, ótica, acústica, química dos produtos naturais, análise de fenômenos físicos e químicos, meteorologia e climatologia, microbiologia, imunologia e parasitologia, ecologia, nutrição, zoologia, dentre outros, considerando o contexto local e as possibilidades de \textbf{oferta} pelos sistemas de ensino ; > > IV. ciências humanas e sociais aplicadas : aprofundamento de conhecimentos \textbf{estruturantes} para aplicação de diferentes conceitos em contextos sociais e de trabalho, estruturando arranjos curriculares que permitam estudos em relações sociais, modelos econômicos, processos políticos, pluralidade cultural, historicidade do universo, do homem e natureza, dentre outros, considerando o contexto local e as possibilidades de \textbf{oferta} pelos sistemas de ensino ; > > V. formação \textbf{técnica} e \textbf{profissional} : desenvolvimento de programas educacionais inovadores e \textbf{atualizados} que promovam efetivamente a \textbf{qualificação} \textbf{profissional} dos estudantes para o mundo do trabalho, objetivando sua \textbf{habilitação} \textbf{profissional} tanto para o desenvolvimento de vida e \textbf{carreira}, quanto para adaptar -se às novas condições ocupacionais e às exigências do mundo do trabalho contemporâneo e suas contínuas transformações, em condições de competitividade, produtividade e inovação, considerando o contexto local e as possibilidades de \textbf{oferta} pelos sistemas de ensino. > > § 1º Os \textbf{itinerários} \textbf{formativos} devem considerar as demandas e necessidades do mundo contemporâneo, estar \textbf{sintonizados} com os diferentes interesses dos estudantes e sua inserção na sociedade, o contexto local e as possibilidades de \textbf{oferta} dos sistemas e instituições de ensino. > > § 2º Os \textbf{itinerários} \textbf{formativos} orientados para o aprofundamento e ampliação das aprendizagens em áreas do conhecimento devem garantir a apropriação de procedimentos cognitivos e uso de metodologias que favoreçam o protagonismo juvenil, e organizar -se em torno de um ou mais dos seguintes eixos \textbf{estruturantes} : > > I. investigação científica : supõe o aprofundamento de conceitos fundantes das ciências para a interpretação de ideias, fenômenos e processos para serem utilizados em procedimentos de investigação voltados ao enfrentamento de situações cotidianas e demandas locais e coletivas, e a proposição de intervenções que considerem o desenvolvimento local e a melhoria da \textbf{qualidade} de vida da comunidade ; > > II. processos criativos : supõe o uso e o aprofundamento do conhecimento científico na construção e criação de experimentos, modelos, protótipos para a criação de processos ou produtos que \textbf{atendam} a demandas pela resolução de problemas identificados na sociedade ; > > III. mediação e intervenção sociocultural : supõe a mobilização de conhecimentos de uma ou mais áreas para mediar conflitos, promover entendimento e implementar soluções para questões e problemas identificados na comunidade ; > > IV. \textbf{empreendedorismo} : supõe a mobilização de conhecimentos de diferentes áreas para a formação de organizações com variadas \textbf{missões} voltadas ao desenvolvimento de produtos ou prestação de serviços inovadores com o uso das tecnologias. > > § 3º \textbf{Itinerários} \textbf{formativos} integrados podem ser \textbf{ofertados} por meio de arranjos curriculares que combinem mais de uma área de conhecimento e da formação \textbf{técnica} e \textbf{profissional}. > > § 4º A definição de \textbf{itinerários} \textbf{formativos} previstos neste artigo e dos seus respectivos arranjos curriculares deve ser orientada pelo \textbf{perfil} de saída almejado para o estudante com base nos Referenciais para a Elaboração dos \textbf{Itinerários} \textbf{Formativos}, e deve ser estabelecido pela instituição ou \textbf{rede} de ensino, considerando os interesses dos estudantes, suas perspectivas de continuidade de estudos no nível pós-secundário e de inserção no mundo do trabalho. > > § 5º Os \textbf{itinerários} \textbf{formativos} podem ser organizados por meio da \textbf{oferta} de diferentes arranjos curriculares, dada a relevância para o contexto local e a possibilidade dos sistemas de ensino. > > § 6º Os sistemas de ensino devem garantir a \textbf{oferta} de mais de um \textbf{itinerário} \textbf{formativo} em cada \textbf{município}, em áreas distintas, permitindo -lhes a escolha, dentre diferentes arranjos curriculares, \textbf{atendendo} assim a heterogeneidade e pluralidade de condições, interesses e aspirações. > > § 7º A critério dos sistemas de ensino, os \textbf{currículos} do ensino médio podem considerar competências \textbf{eletivas} complementares do estudante como forma de ampliação da \textbf{carga} \textbf{horária} do \textbf{itinerário} \textbf{formativo} escolhido, \textbf{atendendo} ao projeto de vida do estudante. > > § 8º A \textbf{oferta} de \textbf{itinerários} \textbf{formativos} deve considerar as possibilidades estruturais e de recursos das instituições ou \textbf{redes} de ensino. > > § 9º Para garantir a \textbf{oferta} de diferentes \textbf{itinerários} \textbf{formativos}, podem ser estabelecidas parcerias entre diferentes instituições de ensino, desde que sejam previamente credenciadas pelos sistemas de ensino, podendo os \textbf{órgãos} normativos em conjunto atuarem como harmonizador dos critérios para credenciamento. > > § 10.
Os sistemas de ensino devem estabelecer o regramento do processo de escolha do \textbf{itinerário} \textbf{formativo} pelo estudante. > > § 11.
As instituições ou \textbf{redes} de ensino devem orientar os estudantes no processo de escolha do seu \textbf{itinerário} \textbf{formativo}. > > § 12.
O estudante pode mudar sua escolha de \textbf{itinerário} \textbf{formativo} ao longo de seu \textbf{curso}, desde que : > > I. resguardadas as possibilidades de \textbf{oferta} das instituições ou \textbf{redes} de ensino ; > II. respeitado o instrumento normativo específico do sistema de ensino. > > § 13.
Os sistemas de ensino devem garantir formas de aproveitamento de estudos realizados com êxito para o estudante em processo de transferência entre instituições ou \textbf{redes} de ensino ou em caso de mudança de \textbf{itinerário} \textbf{formativo} ao longo de seu \textbf{curso}. > > § 14.
O \textbf{itinerário} \textbf{formativo} na formação \textbf{técnica} \textbf{profissional} deve observar a \textbf{integralidade} de ocupações \textbf{técnicas} reconhecidas pelo setor produtivo, tendo como referência a Classificação Brasileira de Ocupações ( CBO ). ( BRASIL, 2018 ) Ainda em dezembro de 2018, o MEC publicou a \textbf{Portaria} nº 1.432, de 28 de dezembro de 2018, que estabeleceu os referenciais para elaboração dos \textbf{itinerários} \textbf{formativos}, conforme preveem as \textbf{Diretrizes} Curriculares Nacionais do Ensino Médio.
Além dos normativos e \textbf{legislação} \textbf{supramencionados}, que vêm orientando os Estados na elaboração dos seus referenciais curriculares para a etapa do Ensino Médio, outras \textbf{políticas} públicas foram implementadas no país, para impulsionar a \textbf{reforma} do Ensino Médio.
Foram elas : - \textbf{Portaria} MEC nº 649, de 10 de julho de 2018, que \textbf{institui} o Programa de Apoio ao Novo Ensino Médio ; - \textbf{Portaria} MEC nº 1.023, de 4 de outubro de 2018, que estabelece \textbf{diretrizes}, parâmetros e critérios para a realização de avaliação de impacto do Programa de Fomento às Escolas de Ensino Médio em Tempo Integral ( EMTI ) e seleção de novas unidades escolares para o Programa ; - \textbf{Portaria} MEC nº 1.024, de 4 de outubro de 2018, que define as \textbf{diretrizes} do apoio financeiro por meio do Programa Dinheiro Direto na Escola às unidades escolares pertencentes às secretarias participantes do Programa de Apoio ao Novo Ensino Médio.
O Conselho Estadual da Educação da Bahia, em 17 de dezembro de 2019, publicou a Resolução CEE nº 137, que > fixa normas complementares para a implementação da Base Nacional Comum Curricular – BNCC, nas \textbf{redes} de ensino e nas instituições escolares integrantes dos sistemas de ensino, na Educação Básica do Estado da Bahia e dá outras \textbf{providências} ( BAHIA, 2019 ).
Dentre as normas complementares para a implementação da Base Nacional Comum Curricular – BNCC, no Estado da Bahia, destaca -se o Artigo 25 que, > No atendimento à parte diversificada, no que tange ao complemento \textbf{previsto} no § 1º do Art. 35 -A da LDB, incluem -se as temáticas seguintes, recomendando -se às instituições escolares a inserção de unidades de ensino conexas aos assuntos na programação curricular : > > I.
Abordagem Territorial como uma \textbf{política} de Estado, seus principais instrumentos ( Lei nº 13.214, de 29 de dezembro de 2014 ) e seus Planos Territoriais de Desenvolvimento Sustentável ( PTDS ), com ênfase na participação social e governança territorial como práticas cidadãs para o desenvolvimento sustentável, inclusivo e colaborativo ; > > II.
Gestão territorial, interfaces com a agenda da sociobiodiversidade e da agroecologia : arranjos de desenvolvimento local e das cadeias produtivas, inclusão produtiva de povos/comunidades tradicionais e estímulo ao \textbf{fortalecimento} das estratégias do desenvolvimento rural ; > > III.
Corredores Ecológicos nos Territórios de Identidade à luz da ecologia da paisagem : planejamento de turismo local, sua institucionalização, sociobiodiversidade e práticas de observação de paisagens, de grutas, de árvores, \textbf{cursos} e espelhos d’água – onde existirem, de aves e outros animais silvestres de pequeno porte ; > > IV.
Cidades e aglomerados populacionais : o paradigma do planejamento ambiental e da ecologia da paisagem, sociobiodiversidade e integrações entre sistemas ecológicos, relações cidade e campo e o contexto das articulações metrópole-região, lógicas de povoamento ante a expansão do desenvolvimento socioeconômico e os modais de transportes na logística do desenvolvimento regional ; > > V.
Bacias hidrográficas da Bahia : biomas, importância bio-socio-ambiental, vetores \textbf{estruturantes} da dimensão sócio-econômica, contribuição sócio-histórica e econômica e culturas ribeirinhas, gestão das águas – comitês de bacias e sua lógica de funcionamento ; > > VI.
Regiões biogeográficas na Bahia : paisagens, ecossistemas, proteção, corredores ecológicos, uso sustentável/comunidades sustentáveis, serviços ecossistêmicos, estudos de priorizações, índices de risco ecológico e cumprimentos de metas de conservação ; > > VII.
Territórios e Etnias : Espaços Quilombolas – marcas da ancestralidade e do senso de pertencimento : diacríticos para a (re)construção identitária.
O lugar da educação para as relações etnicorraciais, da Lei nº 10.639 de 9 de janeiro de 2003 ; > > VIII.
Territórios e Etnias : Espaços Indígenas – direitos territoriais, lutas e resistência ; etnografia e heranças histórico-culturais ; etnodesenvolvimento como \textbf{perfil} de projetos de futuro formulados pelos povos indígenas.
O lugar da educação para as relações etnicorraciais, da Lei nº 11.645 de 10 de março de 2008 ; > > IX.
Territórios, Memórias e Pertencimentos : os movimentos sociais populares – a ruptura com o poder colonial e a utopia de um governo com igualdade racial ( Revoltados Búzios ) ; o 2 de julho no contexto da consolidação da independência política do Brasil ; a saga heroica no sertão de Canudos e a representação do diálogo entre histórias, memórias e identidades da história nacional, regional e local nas diferentes temporalidades. > > X.
Educação em Práticas Corporais : as diferentes manifestações da cultura lúdica dos territórios de identidade do Estado da Bahia e suas expressões, principalmente aquelas de origem de matriz afro-brasileira e indígena. > > XI.
Territórios da Bahia, variações linguísticas e interculturalidades : combinação de traços culturais e a singularização de sujeitos – regiões, linguagem como atividade social, processos linguísticos dos falares baianos, cultura de linguagem e estratégias para o tratamento da variação linguística nas escolas. > > § 1º A inclusão dessas temáticas \textbf{demarca} um conjunto de aspectos importantes à delimitação de fatos representativos ao contexto situacional do Estado, assinalando -se que as instituições escolares podem apresentar temáticas outras, sinalizadas pelas propostas pedagógicas aprovadas pelos seus \textbf{órgãos} competentes. > > § 2º Em obediência ao disposto no § 1º do Art. 35 -A da LDB, acentua -se que essas recomendações traduzem a especificidade da disposição legal quanto às características regionais/territoriais e locais, envolvendo aspectos históricos, culturais, econômicos e ambientais.
Os marcos legais e normativos destacados neste \textbf{capítulo} tiveram como objetivo apresentar as \textbf{diretrizes} que organizam a educação brasileira e baiana, incluindo as normativas atuais que \textbf{amparam} a \textbf{política} curricular do Ensino Médio, evidenciando quais são centrais para a organização do \textbf{currículo} do Ensino Médio baiano.
Versa, ainda, informar aos/às leitores/as que o Volume 3 do DCRB terá como foco a organização curricular das Modalidades de Ensino, contudo se faz necessário considerar nesta versão do documento os normativos, nacionais e \textbf{estaduais}, que orientam tais modalidades, compreendendo que os/as estudantes com deficiência, indígenas, quilombolas, do campo, jovens e adultos, do Ensino Médio, não necessariamente estão \textbf{cursando} nessas modalidades de ensino, poderão \textbf{cursar} a etapa do Ensino Médio em uma das \textbf{ofertas} apresentadas neste documento.
A seguir serão apresentadas duas tabelas, uma que apresentará os normativos que regem as Modalidades de Ensino e outra com os normativos que embasam os Temas Integradores do Documento Curricular Referencial da Bahia.

% matched lemmas: adolescente, amparar, atender, atualizar, capítulo, carga, carreira, currículo, cursar, curso, demarcar, destinar, diretriz, duzentos, efetivação, eletivo, empreendedorismo, entidade, equidade, estadual, estruturante, flexível, formativo, fortalecimento, garantia, habilitação, horário, instituir, integralidade, itinerário, legislação, missão, municipal, município, normatizar, oferta, ofertar, perfil, política, portaria, precípuo, previsto, profissional, prosseguimento, providência, qualidade, qualificação, rede, reforma, responsabilizar, revogar, sintonizar, socioeducacional, supramencionado, trabalhador, técnico, viabilizar, órgão
\end{textsample}
